\documentclass[../../../include/open-logic-section]{subfiles}

\begin{document}

\olfileid[hi]{sth}{card-arithmetic}{fix}
\olsection{$\aleph$-नियत बिंदु}

\olref[spine][]{chap} में हमने सुझाया था कि प्रतिस्थापन के लिए औचित्य की
आवश्यकता है, क्योंकि वह पदानुक्रम को बहुत ऊँचा होने के लिए बाध्य करता है।
अब कुछ कार्डिनल अंकगणित कर लेने के बाद हम पदानुक्रम की ऊँचाई का एक छोटा
चित्रण दे सकते हैं।

स्पष्टतः $0<\aleph_0$, और $1<\aleph_1$, और $2<\aleph_2$\ldots और वास्तव
में आकार का अंतर प्रत्येक चरण के साथ केवल \emph{बढ़ता} जाता है। अतः यह
अनुमान लगाना आकर्षक है कि प्रत्येक क्रमसूचक $\kappa$ के लिए
$\kappa<\aleph_\kappa$।

किंतु $\ZFC$ दिए होने पर यह अनुमान \emph{असत्य} है। वास्तव में हम सिद्ध कर
सकते हैं कि \emph{$\aleph$-नियत बिंदु} हैं, अर्थात् ऐसे कार्डिनल $\kappa$
जिनके लिए $\kappa=\aleph_\kappa$।

\begin{prop}\ollabel{alephfixed}
कोई $\aleph$-नियत बिंदु है।
\end{prop}

\begin{proof}
पुनरावर्तन का उपयोग करके परिभाषित करें:
\begin{align*}
	\kappa_0 &= 0\\
	\kappa_{n+1} &= \aleph_{\kappa_n}\\
	\kappa&= \bigcup_{n < \omega}\kappa_n
\end{align*}
अब \olref[cardinals][classing]{unioncardinalscardinal} से $\kappa$ एक
कार्डिनल है। किंतु अब:
\[
	\kappa= \bigcup_{n < \omega} \kappa_{n+1} = 
	\bigcup_{n < \omega}\aleph_{\kappa_n} = 
	\bigcup_{\alpha < \kappa}\aleph_\alpha = \aleph_\kappa
\]
%	By construction, $\kappa$ is the least cardinal greater than each
%	$\kappa_n$. So $\aleph_\kappa$ is the least cardinal greater than
%	each $\aleph_{\kappa_n}$, and hence greater than each
%	$\kappa_{n+1}$. But equally $\kappa$ is the least cardinal greater
%	than each $\kappa_{n+1} = \aleph_{\kappa_n}$. So $\kappa=
%	\aleph_\kappa$.
\end{proof}

बूलोस ने एक बार ठीक उसी $\aleph$-नियत बिंदु के बारे में लेख लिखा जिसे हमने
अभी निर्मित किया है। अपने लेख के आरंभ में $\kappa$ के अस्तित्व पर ध्यान
दिलाने के बाद उन्होंने कहा:
\begin{quote}
	[$\kappa$] उन लोगों की दृष्टि में, जिनका समुच्चय सिद्धांत से पहले कोई
	परिचय नहीं रहा है, एक \emph{बहुत बड़ी} संख्या [है]---मेरे विचार में इतनी
	बड़ी कि वह किसी भी ऐसे सिद्धांत की सत्यता पर प्रश्न उठाती है जिसका एक
	कथन यह दावा हो कि कम-से-कम $\kappa$ वस्तुएँ हैं।
	\citep[p.~257]{Boolos2000}
\end{quote}
और अंततः उन्होंने अपना लेख इस प्रश्न के साथ समाप्त किया:
\begin{quote}
	[क्या] हमें संदेह है कि कहानी के आरंभ में स्थिति चाहे जैसी भी रही हो,
	यहाँ तक पहुँचते-पहुँचते पहिए व्यर्थ घूम रहे हैं और हम अब किसी वास्तविक
	स्थिति का वर्णन नहीं सुन रहे हैं?
	\citep[p.~268]{Boolos2000}
\end{quote}
यदि हम वास्तव में ``किसी वास्तविक स्थिति'' से आगे निकल गए हैं, तो दोष सीधे
प्रतिस्थापन पर लगाना होगा। क्योंकि यही वह अभिगृहीत है जो हमारे प्रमाण को
काम करने देता है। ऐसी स्थिति में, संभवतः बूलोस को कुछ दशक पहले किया अपना
यह दावा फिर देखना होगा कि प्रतिस्थापन के ``कोई अवांछनीय'' परिणाम नहीं हैं
(\olref[replacement][extrinsic]{sec} देखें)।

किंतु क्या $\kappa$ का अस्तित्व इतना बुरा है? यहाँ रसेल के
\emph{ट्रिस्ट्रम शैंडी विरोधाभास} पर विचार करना सहायक हो सकता है। ट्रिस्ट्रम
शैंडी अपने जीवन का विवरण डायरी में लिखता है, किंतु एक दिन दर्ज करने में उसे
एक वर्ष लगता है। प्रत्येक बीतते वर्ष के साथ ट्रिस्ट्रम और पीछे रह जाता है:
एक वर्ष बाद उसने केवल एक दिन दर्ज किया और 364 दिन बिना दर्ज किए जी चुका है;
दो वर्ष बाद उसने केवल दो दिन दर्ज किए और 728 दिन बिना दर्ज किए जी चुका है;
तीन वर्ष बाद उसने केवल तीन दिन दर्ज किए और 1092 दिन बिना दर्ज किए जी चुका
है\dots\footnote{अधिवर्षों को भूलते हुए।} फिर भी यदि ट्रिस्ट्रम
\emph{अमर} है, तो वह प्रत्येक दिन दर्ज कर लेगा, क्योंकि वह अपने जीवन के
$n$वें वर्ष में $n$वाँ दिन दर्ज करेगा। इसलिए ``समय के अंत में'' ट्रिस्ट्रम
के पास पूरी डायरी होगी।

अब: यह उस विचार से इतना भिन्न क्यों है कि $\kappa$ तक $\alpha$,
$\aleph_\alpha$ से छोटा---और वास्तव में निरंतर, अत्यधिक छोटा---है, और
$\kappa$ पर हम बराबरी पा लेते हैं तथा $\kappa=\aleph_\kappa$?

इसे अलग रखते हुए और यह मानते हुए कि हम $\ZFC$ को स्वीकार करते हैं, आइए
नियत-बिंदु निर्माणों के विषय में कुछ और रोचक परिणामों के साथ समाप्त करें।
अगले तीन परिणाम सहज रूप से स्थापित करते हैं कि कोई (अतुच्छ) बिंदु है जहाँ
पदानुक्रम जितना ऊँचा है उतना ही चौड़ा भी है:

\begin{prop}\ollabel{bethfixed}
कोई $\beth$-नियत बिंदु है, अर्थात् कोई ऐसा $\kappa$ कि
$\kappa=\beth_\kappa$।
\end{prop}

\begin{proof}
\olref{alephfixed} की तरह, ``$\aleph$'' के स्थान पर ``$\beth$'' का उपयोग
करके।
\end{proof}

\begin{prop}\ollabel{stagesize}
$\card{V_{\omega+\alpha}}=\beth_{\alpha}$। यदि
$\omega\ordtimes\omega\leq\alpha$, तो $\card{V_\alpha}=\beth_\alpha$।
\end{prop}

\begin{proof}
पहला दावा सरल परिमितेतर आगमन से सिद्ध होता है। दूसरा दावा इसलिए निष्पन्न
होता है कि यदि $\omega\ordtimes\omega\leq\alpha$, तो
$\omega+\alpha=\alpha$। इसे स्थापित करने के लिए हम
\olref[ord-arithmetic][]{chap} के क्रमसूचक अंकगणित संबंधी तथ्यों का उपयोग
करते हैं। पहले ध्यान दें कि
$\omega\ordtimes\omega=\omega\ordtimes(1\ordplus\omega)=
(\omega\ordtimes1)\ordplus(\omega\ordtimes\omega)=
\omega\ordplus(\omega\ordtimes\omega)$। अब यदि
$\omega\ordtimes\omega\leq\alpha$, अर्थात् किसी $\beta$ के लिए
$\alpha=(\omega\ordtimes\omega)\ordplus\beta$, तो
$\omega\ordplus\alpha=\omega\ordplus((\omega\ordtimes\omega)\ordplus
\beta)=(\omega\ordplus(\omega\ordtimes\omega))\ordplus\beta=
(\omega\ordtimes\omega)\ordplus\beta=\alpha$।
\end{proof}

\begin{cor}
कोई ऐसा $\kappa$ है कि $\card{V_\kappa}=\kappa$।
\end{cor}

\begin{proof}
$\kappa$ को \olref{bethfixed} द्वारा दिया कोई $\beth$-नियत बिंदु मानें।
स्पष्टतः $\omega\ordtimes\omega<\kappa$। अतः \olref{stagesize} से
$\card{V_\kappa}=\beth_\kappa=\kappa$।
\end{proof}

$V_\kappa$ के नीचे उतने ही चरण हैं जितने $V_\kappa$ के !!{element}। अतः
सहज रूप से $V_\kappa$ जितना ऊँचा है उतना ही चौड़ा है। यह बहुत
ट्रिस्ट्रम-शैंडी जैसा है: हम \emph{घात समुच्चय} लेकर एक चरण से अगले चरण तक
जाते हैं और इस प्रकार प्रत्येक कदम पर अपने पदानुक्रम को \emph{कहीं} बड़ा
बनाते हैं। किंतु ``अंत में'', अर्थात् चरण $\kappa$ पर, पदानुक्रम की चौड़ाई
उसकी ऊँचाई को पकड़ लेती है।

कोई पूछ सकता है: \emph{पदानुक्रम की चौड़ाई उसकी ऊँचाई के बराबर कितनी बार
होती है?} उत्तर है: \emph{जितनी बार क्रमसूचक हैं।} किंतु इसके लिए कुछ
व्याख्या चाहिए।

हम एक पद $\tau$ इस प्रकार परिभाषित करते हैं। किसी भी $A$ के लिए रखें:
\begin{align*}
	\tau_0(A) & \defis \card{A}\\
	\tau_{n+1}(A) & \defis \beth_{\tau_n(A)}\\
	\tau(A) & \defis \bigcup_{n < \omega}\tau_n(A)
\intertext{\olref{bethfixed} की तरह किसी भी $A$ के लिए $\tau(A)$ एक
$\beth$-नियत बिंदु है, और तुच्छतः $\card{A}<\tau(A)$। अब इस पुनरावर्ती
परिभाषा पर विचार करें:}
	W_0 &\defis 0\\
	W_{\alpha + 1} & \defis \tau(W_\alpha)\\
	W_\alpha & \defis \bigcup_{\beta < \alpha} W_\beta
	\text{, जब $\alpha$ सीमा क्रमसूचक हो}
\end{align*}
यह निर्माण सभी क्रमसूचकों के लिए परिभाषित है। अतः सहज रूप से $W$
क्रमसूचकों से $\beth$-नियत बिंदुओं तक ``!!a{injection}'' है। और ठीक पहले
की तरह किसी भी $\alpha$ के लिए $V_{W_\alpha}$ जितना ऊँचा है उतना ही चौड़ा
है।

\end{document}
